% GENERATED FILE. Edit canonical lessons, then run npm run generate:books.
% canonical-source-hash: fnv1a64:957130e12f222e35
% canonical-lessons: FR-C356-emploi, FR-C356-chomage, FR-C356-chomeur, FR-C356-embaucher, FR-C356-recruter

\chapter{Sans emploi}
\label{ch:fr-sans-emploi}
\hlchaptermodality{356}

\begin{chapteropening}
\textbf{By the end of this chapter:} \emph{I can talk about looking for a job: unemployment, being out of work, and companies that hire and recruit.}

Complete the last lesson of chapter 356: I can talk about looking for a job: unemployment, being out of work, and companies that hire and recruit.
\end{chapteropening}

\section[l'emploi]{l'emploi — the job, employment; the use}

\label{lesson:FR-C356-emploi}



\begin{quote}
Before the new one: say the French for to resign, then the French for the notice.
\end{quote}

\subsection*{You'll want to know: l'emploi}
\textbf{l'emploi} — \textquotedblleft{}the job, employment; the use\textquotedblright{}.

\begin{quote}
Mon frère cherche un emploi dans une banque.
\end{quote}

\begin{quote}
\emph{My brother is looking for a job in a bank.}
\end{quote}

\emph{Un emploi du temps} — a timetable. \emph{Le mode d'emploi} — the instructions. Masculine.

\begin{cousinweb}
\emph{emploi} is formed from \emph{employer} (EtymDB), which is inherited from Latin \emph{implicāre} \textquotedblleft{}to enfold, to involve\textquotedblright{} (Lewis \& Short).
\end{cousinweb}

\begin{grammarlens}[title={What you've built}]
One more word for telling, explaining and describing at B1.
\end{grammarlens}

\subsection*{Your turn}
\begin{itemize}
\raggedright
  \item \emph{Say it:} \emph{l'emploi}
  \item \emph{Say it:} \emph{l'emploi}, once more
  \item \emph{Say it:} a sentence of your own with \emph{l'emploi}
  \item \emph{Recall:} say the French for to resign, then the French for the notice, then say \emph{l'emploi} again
\end{itemize}

\begin{tcolorbox}[breakable,colback=teal!4,colframe=teal!35!black,title={Before you move on}]
What does l'emploi mean? (\textquotedblleft{}The job, employment; the use\textquotedblright{}.) Say it once more.
\end{tcolorbox}
\section[le chômage]{le chômage — the unemployment}

\label{lesson:FR-C356-chomage}



\begin{quote}
Before the new one: say the French for the notice, then the French for the job.
\end{quote}

\subsection*{You'll want to know: le chômage}
\textbf{le chômage} — \textquotedblleft{}the unemployment\textquotedblright{}.

\begin{quote}
Il est au chômage depuis trois mois.
\end{quote}

\begin{quote}
\emph{He has been out of work for three months.}
\end{quote}

\emph{Être au chômage} — to be out of work. \emph{Toucher le chômage} — to get unemployment benefit; in Québec the benefit is \emph{l'assurance-emploi}.

\begin{cousinweb}
\emph{chômage} is built on \emph{chômer} \textquotedblleft{}to stop work\textquotedblright{}. The verb is usually traced to Late Latin \emph{caumāre} \textquotedblleft{}to rest during the heat\textquotedblright{}, from Greek \emph{kaûma} \textquotedblleft{}burning heat\textquotedblright{}.
\end{cousinweb}

\begin{grammarlens}[title={What you've built}]
One more word for telling, explaining and describing at B1.
\end{grammarlens}

\subsection*{Your turn}
\begin{itemize}
\raggedright
  \item \emph{Say it:} \emph{le chômage}
  \item \emph{Say it:} \emph{le chômage}, once more
  \item \emph{Say it:} a sentence of your own with \emph{le chômage}
  \item \emph{Recall:} say the French for the notice, then the French for the job, then say \emph{le chômage} again
\end{itemize}

\begin{tcolorbox}[breakable,colback=teal!4,colframe=teal!35!black,title={Before you move on}]
What does le chômage mean? (\textquotedblleft{}The unemployment\textquotedblright{}.) Say it once more.
\end{tcolorbox}
\section[le chômeur]{le chômeur — the unemployed person}

\label{lesson:FR-C356-chomeur}



\begin{quote}
Before the new one: say the French for the job, then the French for the unemployment.
\end{quote}

\subsection*{You'll want to know: le chômeur}
\textbf{le chômeur} — \textquotedblleft{}the unemployed person\textquotedblright{}.

\begin{quote}
Ce chômeur cherche du travail tous les jours.
\end{quote}

\begin{quote}
\emph{This unemployed man looks for work every day.}
\end{quote}

Feminine \emph{chômeuse}.

\begin{cousinweb}
\emph{chômeur} is built from \emph{chômer} \textquotedblleft{}to stop work\textquotedblright{} and \emph{-eur}. The verb is usually traced to Late Latin \emph{caumāre} \textquotedblleft{}to rest during the heat\textquotedblright{}, from Greek \emph{kaûma} \textquotedblleft{}burning heat\textquotedblright{}.
\end{cousinweb}

\begin{grammarlens}[title={What you've built}]
One more word for telling, explaining and describing at B1.
\end{grammarlens}

\subsection*{Your turn}
\begin{itemize}
\raggedright
  \item \emph{Say it:} \emph{le chômeur}
  \item \emph{Say it:} \emph{le chômeur}, once more
  \item \emph{Say it:} a sentence of your own with \emph{le chômeur}
  \item \emph{Recall:} say the French for the job, then the French for the unemployment, then say \emph{le chômeur} again
\end{itemize}

\begin{tcolorbox}[breakable,colback=teal!4,colframe=teal!35!black,title={Before you move on}]
What does le chômeur mean? (\textquotedblleft{}The unemployed person\textquotedblright{}.) Say it once more.
\end{tcolorbox}
\section[embaucher]{embaucher — to hire, to take on}

\label{lesson:FR-C356-embaucher}



\begin{quote}
Before the new one: say the French for the unemployment, then the French for the unemployed person.
\end{quote}

\subsection*{You'll want to know: embaucher}
\textbf{embaucher} — \textquotedblleft{}to hire, to take on\textquotedblright{}.

\begin{quote}
Le restaurant embauche des serveurs pour l'été.
\end{quote}

\begin{quote}
\emph{The restaurant is taking on waiters for the summer.}
\end{quote}

\emph{Le restaurant embauche} — the restaurant is hiring.

\begin{cousinweb}
\emph{embaucher} is usually explained as built from \emph{en-} and an old word \emph{bauche} \textquotedblleft{}workshop\textquotedblright{}: to bring someone into the workshop. The origin of \emph{bauche} is not settled.
\end{cousinweb}

\begin{grammarlens}[title={What you've built}]
One more word for telling, explaining and describing at B1.
\end{grammarlens}

\subsection*{Your turn}
\begin{itemize}
\raggedright
  \item \emph{Say it:} \emph{embaucher}
  \item \emph{Say it:} \emph{embaucher}, once more
  \item \emph{Say it:} a sentence of your own with \emph{embaucher}
  \item \emph{Recall:} say the French for the unemployment, then the French for the unemployed person, then say \emph{embaucher} again
\end{itemize}

\begin{tcolorbox}[breakable,colback=teal!4,colframe=teal!35!black,title={Before you move on}]
What does embaucher mean? (\textquotedblleft{}To hire, to take on\textquotedblright{}.) Say it once more.
\end{tcolorbox}
\section[recruter]{recruter — to recruit}

\label{lesson:FR-C356-recruter}



\begin{quote}
Before the new one: say the French for the unemployed person, then the French for to hire.
\end{quote}

\subsection*{You'll want to know: recruter}
\textbf{recruter} — \textquotedblleft{}to recruit\textquotedblright{}.

\begin{quote}
L'entreprise recrute des ingénieurs.
\end{quote}

\begin{quote}
\emph{The company is recruiting engineers.}
\end{quote}

The person recruited is \emph{une recrue}.

\begin{cousinweb}
\emph{recruter} is built on \emph{recrue} \textquotedblleft{}a recruit\textquotedblright{} (EtymDB), the old participle of \emph{recroître} \textquotedblleft{}to grow again\textquotedblright{}: new growth for a group. \emph{croître} is inherited from Latin \emph{crēscere} \textquotedblleft{}to grow\textquotedblright{} (Lewis \& Short).
\end{cousinweb}

\begin{grammarlens}[title={What you've built}]
One more word for telling, explaining and describing at B1.
\end{grammarlens}

\subsection*{Your turn}
\begin{itemize}
\raggedright
  \item \emph{Say it:} \emph{recruter}
  \item \emph{Say it:} \emph{recruter}, once more
  \item \emph{Say it:} a sentence of your own with \emph{recruter}
  \item \emph{Recall:} say the French for the unemployed person, then the French for to hire, then say \emph{recruter} again
\end{itemize}

\begin{tcolorbox}[breakable,colback=teal!4,colframe=teal!35!black,title={Before you move on}]
What does recruter mean? (\textquotedblleft{}To recruit\textquotedblright{}.) Say it once more.
\end{tcolorbox}
